% This file was created with the script capacityfactor.py
\mode* % only put content inside frame environments on slides

\begin{table}\centering
    \caption{Example of a table generated in Python}\label{tab:pythontable}
    \begin{tabular}{lrl}  
        \toprule
        parameter               &  value            & unit              \\[1ex]
        \textit{inputs:}        &                   &                   \\
        nominal power           &  \num{39}       & \unit{kW}         \\
        electricity produced    &  \num{90.4}       & \unit{MWh/year}   \\[1ex]
        \textit{calculated:}    &                   &                   \\
        capacity factor         &  \num{26.5}       & \unit{\percent}   \\
        equiv. full load hours  &  \num{2318}       & \unit{\h}         \\
        \bottomrule
    \end{tabular}
\end{table}

\begin{frame}<presentation>{Example of a slide generated in Python (with script capacityfactor.py)}
An engine with a nominal power of \qty{39}{\kW} produces \qty{90.4}{\MWh} of electricity in a year.
Calculate the capacity factor (CF) and the number of equivalent full load hours (EFLH).

\begin{align*}
    \intertext{The capacity factor}
    &\up{CF} = \frac{\qty{90}{\MWh\per\yr}}{\qty{39}{\kW}\cdot\qty{8760}{\h\per\yr}}=\qty{26.5}{\percent} 
    \intertext{Equivalent hours of full load operation}
    &\up{EFLH} = \frac{\qty{90}{\MWh\per\yr}}{\qty{39}{\kW}}=\qty{2318}{\h} 
\end{align*}
\end{frame}
